%% =====================================================================
%% Multimodal Emotion-Cause Pair Extraction with Necessity Evidence,
%% LLM-Distilled Counterfactual Reasoning, and Complementarity-Aware Fusion
%%
%% Draft prepared in the Elsevier elsarticle format, suitable for
%% submission to CCF-C international journals such as Neurocomputing,
%% Knowledge-Based Systems, or Expert Systems with Applications.
%%
%% NOTE (draft): all reported numbers in Section "Experiments" are
%% PLACEHOLDER values to be replaced by the real experimental results.
%% Search for "TODO" to locate every number that must be updated.
%% =====================================================================
\documentclass[preprint,authoryear,12pt]{elsarticle}

\usepackage{amsmath,amssymb,amsfonts}
\usepackage{graphicx}
\usepackage{booktabs}
\usepackage{multirow}
\usepackage{array}
\usepackage{url}
\usepackage[hidelinks]{hyperref}
\usepackage{xcolor}

\newcommand{\todo}[1]{\textcolor{red}{[TODO: #1]}}
\newcommand{\Real}{\mathbb{R}}
\newcommand{\sg}{\mathrm{sg}}
\newcommand{\Enc}{\mathrm{Enc}}
\newcommand{\Fuse}{\mathrm{Fuse}}

\journal{Neurocomputing}

\begin{document}

\begin{frontmatter}

\title{Necessity, Reasoning, and Position: Complementarity-Aware Evidence
Fusion for Multimodal Emotion-Cause Pair Extraction in Conversations}

\author[a]{First Author}
\ead{first.author@example.edu}
\author[a]{Second Author}
\author[a]{Corresponding Author\corref{cor1}}
\ead{corresponding.author@example.edu}
\cortext[cor1]{Corresponding author.}
\affiliation[a]{organization={Department of Computer Science, Shantou University},
  city={Shantou},
  postcode={515063},
  country={China}}

\begin{abstract}
Multimodal emotion-cause pair extraction (MECPE) aims to identify, for every
utterance that expresses an emotion in a conversation, the utterances that cause
that emotion, using textual, acoustic, and visual signals together with speaker
identity. Existing approaches model inter-utterance interaction but rarely ask
\emph{why} a candidate is predicted to be a cause: their decisions are not
grounded in any explicit notion of necessity, and the strong positional regularity
of causes is either ignored or allowed to dominate. We propose a model that
couples a relation-aware multimodal encoder with three complementary sources of
decision evidence: (i) perturbation-based \emph{necessity evidence} that probes how
much each candidate cause and each non-verbal modality contributes to a pair
prediction; (ii) a \emph{bounded positional prior} that exploits the locality of
causes without letting position override semantics; and (iii) \emph{reasoning
evidence} distilled offline from a large language model (LLM) as text-grounded
counterfactual judgements. A complementarity-aware gating network fuses these
signals conditioned on model confidence, teacher reliability, and their agreement,
so that each source compensates the other's weakness; the LLM is queried only
offline and never at inference. Experiments on the benchmark Emotion-Cause-in-Friends
(ECF) corpus show that the proposed method improves the emotion-cause pair $F_1$ over
strong baselines, and ablations confirm that every component contributes. % TODO: tighten claims after final numbers
\end{abstract}

\begin{keyword}
Emotion-cause pair extraction \sep Multimodal conversation understanding \sep
Graph neural networks \sep Counterfactual reasoning \sep Knowledge distillation \sep
Affective computing
\end{keyword}

\end{frontmatter}

% =====================================================================
\section{Introduction}
\label{sec:intro}

Understanding \emph{why} a speaker feels a certain way is central to affective
computing. Beyond recognizing that an utterance conveys an emotion, many
applications---empathetic dialogue systems, opinion mining, and mental-health
support---require identifying the events or utterances that \emph{trigger} that
emotion. The task of emotion-cause pair extraction (ECPE) formalizes this goal:
given a document, jointly extract emotion utterances and their corresponding cause
utterances as pairs \citep{xia2019ecpe}. In conversations, emotions and their
causes are conveyed not only through text but also through tone of voice and facial
expression, which motivates the multimodal emotion-cause pair extraction (MECPE)
task \citep{wang2023mecpe}: for each emotion-bearing utterance in a multi-party
conversation, find the utterances that cause it, using text, audio, and video.

MECPE is challenging for three intertwined reasons. \emph{First}, the label space
is quadratic in the number of utterances and highly imbalanced: only a small
fraction of utterance pairs are genuine emotion-cause pairs. \emph{Second}, a
large share of causes are \emph{self-contained}---the emotion is triggered by the
event described \emph{within} the emotion utterance itself, so the cause coincides
with the emotion utterance. \emph{Third}, causes exhibit strong but not absolute
positional regularity: they tend to lie close to the emotion utterance, which is
informative yet dangerous to over-trust.

Recent models substantially improve MECPE by modeling rich inter-utterance
interaction with graph or transformer architectures and by injecting label
constraints \citep{li2024hilo}. However, two limitations remain. (1) Their pair
decisions are not grounded in any explicit notion of \emph{necessity}: the model
predicts a pair without ever asking whether removing the candidate cause would
actually change the prediction, which makes decisions hard to calibrate and to
explain. (2) External world knowledge---the kind of commonsense causal reasoning
at which large language models (LLMs) excel---is not exploited, even though it is
exactly the missing ingredient for distinguishing genuine triggers from mere
topical co-occurrence. A naive remedy, querying an LLM per pair at inference, is
impractical at the quadratic scale of MECPE and cannot perceive the non-verbal
modalities.

We address these limitations with a model that augments a relation-aware
multimodal encoder with \emph{three complementary sources of decision evidence}.
\textbf{(i) Necessity evidence} (Section~\ref{sec:necessity}) is a model-internal,
perturbation-based signal: we replace the candidate cause---and, separately, each
non-verbal modality of the cause utterance---with a learned, data-anchored neutral
baseline, re-score the pair, and measure the resulting drop. A direct calibration
loss anchors this signal to the pair labels. \textbf{(ii) A bounded positional
prior} (Section~\ref{sec:pos}) encodes relative position as a capped additive term
on the pair score, with a regularizer that prevents position from overriding the
semantic evidence. \textbf{(iii) Reasoning evidence} (Section~\ref{sec:llm}) is
distilled offline from an LLM: for hard pairs, a structured prompt elicits
text-grounded \emph{counterfactual} judgements (necessity, sufficiency, direction,
spuriousness) together with a self-agreement reliability estimate; a lightweight
student then predicts these scalars from the pair representation alone, so that no
LLM is queried at inference.

The three signals are reconciled by a \textbf{complementarity-aware gating network}
(Section~\ref{sec:fusion}) that conditions on the model's confidence, the teacher's
predicted reliability, and the agreement between the model's own necessity signal
and the LLM's necessity judgement. The design realizes an explicit division of
labor: the LLM \emph{compensates} the small model where the latter is uncertain,
while the small model \emph{leads} where the LLM is unreliable or disagrees. We
further introduce a necessity-modulated agreement term so that consensus boosts the
score only when both sources agree a pair \emph{is} causal, preventing
false-positive inflation when both agree it is \emph{not}.

Our contributions are summarized as follows.
\begin{itemize}
  \item We introduce a unified MECPE model that grounds pair decisions in explicit
  \emph{necessity evidence}, obtained by faithful, pair-local perturbations of the
  cause utterance and of each non-verbal modality, and calibrated directly against
  the pair labels.
  \item We distill commonsense \emph{counterfactual reasoning} from an LLM into a
  lightweight student with a dedicated treatment of self-contained causes, so that
  semantically informed external knowledge is available at inference \emph{without}
  any LLM call.
  \item We propose a \emph{complementarity-aware fusion} mechanism, together with a
  bounded positional prior and a necessity-modulated agreement term, that lets each
  evidence source compensate the others' weaknesses under explicit reliability and
  agreement conditioning.
  \item On the benchmark ECF corpus the method improves emotion-cause pair $F_1$
  over strong baselines, and extensive ablations isolate the contribution of each
  component. % TODO: finalize claims after experiments
\end{itemize}

% =====================================================================
\section{Related Work}
\label{sec:related}

\paragraph{Emotion cause analysis}
Emotion cause extraction was first framed as detecting the clauses that cause a
given emotion \citep{lee2010emotion}. \citet{xia2019ecpe} introduced emotion-cause
pair extraction (ECPE), which removes the need for emotion annotations at test time
by jointly extracting emotion and cause clauses as pairs. Subsequent work improved
ECPE with two-dimensional representations \citep{ding2020ecpe2d} and effective
inter-clause ranking \citep{wei2020rankcp}. These methods operate on text alone and
on monologue-style documents.

\paragraph{Multimodal emotion-cause pair extraction}
\citet{wang2023mecpe} extended the task to conversations and to multiple modalities,
releasing the Emotion-Cause-in-Friends (ECF) corpus and establishing the MECPE task.
\citet{li2024hilo} proposed a holistic-interaction model with label constraints that
substantially advanced the state of the art. Our work builds on this line: we adopt
a relation-aware multimodal encoder over the conversation graph, but depart from
prior work by grounding pair decisions in explicit necessity evidence and by
distilling external counterfactual reasoning, rather than relying on interaction
modeling alone.

\paragraph{Graph and multimodal modeling of conversations}
Graph neural networks are a natural fit for conversations, where utterances relate
through speaker identity and reply structure \citep{ghosal2019dialoguegcn}. We use
relation-specific attention in the spirit of relational GCNs \citep{schlichtkrull2018rgcn}
and graph attention networks \citep{velickovic2018gat}. For multimodal fusion, prior
work disentangles modality-invariant and modality-specific factors \citep{hazarika2020misa}
and aligns modalities with cross-modal transformers \citep{tsai2019mult}; we keep the
fusion light and instead focus the modeling effort on causal-style evidence.

\paragraph{Interpretability and counterfactual reasoning}
Perturbation-based attribution explains predictions by measuring sensitivity to
input changes \citep{ribeiro2016lime,lundberg2017shap}. We repurpose this idea
\emph{inside} training: necessity evidence is a model-level signal that both
regularizes and explains pair decisions, distinct from asserting causal relations in
the world. Counterfactual reasoning has deep roots in causal inference
\citep{pearl2009causality}; we obtain weak, text-grounded counterfactual judgements
by prompting an LLM, which is complementary to---not a substitute for---the
model-internal perturbation signal.

\paragraph{LLMs and knowledge distillation}
LLMs exhibit strong commonsense reasoning when prompted to reason step by step
\citep{wei2022cot}, but querying them at the quadratic scale of MECPE, and for
non-verbal signals, is impractical. We therefore distill their reasoning offline into
a small student \citep{hinton2015distilling}, restricting the LLM to text-grounded
judgements and obtaining modality evidence from perturbations instead.

% =====================================================================
\section{Method}
\label{sec:method}

We address MECPE in conversations, where the goal is to identify, for each utterance
that expresses emotion, the utterances that cause it. Our model couples a
relation-aware multimodal encoder with three complementary sources of decision
evidence: necessity evidence (Section~\ref{sec:necessity}), a bounded positional prior
(Section~\ref{sec:pos}), and reasoning evidence distilled offline from an LLM
(Section~\ref{sec:llm}). The signals are combined by a gated fusion module
(Section~\ref{sec:fusion}) and trained jointly (Section~\ref{sec:training}).
Figure~\ref{fig:overview} gives an overview. % TODO: add overview figure (Graph/结.png can be adapted)

\begin{figure}[t]
  \centering
  % TODO: replace with the real architecture figure
  \fbox{\parbox[c][4.5cm][c]{0.9\linewidth}{\centering\itshape
    Placeholder for the model overview figure: relation-aware encoder
    $\rightarrow$ soft candidate gate $\rightarrow$ \{necessity evidence,
    positional prior, distilled reasoning evidence\} $\rightarrow$
    complementarity-aware fusion $\rightarrow$ pair probability.}}
  \caption{Overview of the proposed model. The LLM is used only offline to produce
  distillation targets; at inference the student supplies the reasoning evidence.}
  \label{fig:overview}
\end{figure}

\subsection{Task Formulation}
\label{sec:task}
A conversation is a sequence of $N$ utterances $\mathcal{D}=\{u_1,\ldots,u_N\}$.
Each utterance carries textual, acoustic, and visual features together with a speaker
identity, $u_i=\{x_i^{t},x_i^{a},x_i^{v},s_i\}$. For an ordered pair $(u_i,u_j)$ in
which $u_i$ is the candidate \emph{emotion} utterance and $u_j$ the candidate
\emph{cause} utterance, the model predicts
\begin{equation}
p_{ij}=P\!\left(y_{ij}=1 \mid u_i,u_j,\mathcal{D}\right),
\end{equation}
where $y_{ij}=1$ iff $u_j$ is a cause of the emotion expressed in $u_i$. Two
auxiliary utterance-level heads predict an emotion probability $\hat p^{\,e}_i$ and a
cause probability $\hat p^{\,c}_j$; these supply supervision and drive the candidate
gating of Section~\ref{sec:gate}.

\subsection{Relation-Aware Multimodal Encoder}
\label{sec:encoder}
\paragraph{Multimodal fusion} Each utterance is projected into a shared
$H$-dimensional space by fusing its modalities and speaker identity,
\begin{equation}
h_i^{0}=\Fuse\!\left(x_i^{t},x_i^{a},x_i^{v},s_i\right)\in\Real^{H}.
\end{equation}

\paragraph{Conversation graph} We build a heterogeneous graph over utterances with
four relation types
$\mathcal{R}=\{\textit{self-loop},\textit{same-speaker},\textit{reply/adjacent},
\textit{emotion-transition}\}$: \emph{self-loop} ($i=j$); \emph{same-speaker}
($s_i=s_j,\,i\neq j$); \emph{reply/adjacent} ($|i-j|=1$ or $u_j$ is the reply target
of $u_i$, with a fallback to the nearest preceding cross-speaker utterance); and
\emph{emotion-transition}, where structurally connected turns whose emotion labels
differ are re-typed, taking priority over the structural off-diagonal types. The
emotion-transition relation depends on emotion labels, which are gold during training
and predicted at inference. To avoid a train--inference mismatch, the source label of
each such edge is drawn by scheduled sampling \citep{bengio2015scheduled},
\begin{equation}
\ell_i=
\begin{cases}
\ell_i^{\mathrm{gold}}, & \text{with probability } \varepsilon_t,\\
\ell_i^{\mathrm{pred}}, & \text{with probability } 1-\varepsilon_t,
\end{cases}
\qquad \varepsilon_t:\,1\rightarrow 0,
\end{equation}
where $\ell_i^{\mathrm{pred}}$ is the detached arg-max of a preliminary emotion head.
As $\varepsilon_t$ anneals to $0$, the training-time graph converges to the
inference-time graph. This relation and the evidence modules of
Sections~\ref{sec:necessity} and \ref{sec:llm} are activated only after a warmup of
$K$ epochs; during warmup the graph uses only the three structural relation types.

\paragraph{Relation-aware propagation} Messages are aggregated with relation-specific
attention,
\begin{equation}
\alpha_{ij}^{r}=\operatorname*{softmax}_{j}\!\Big(\mathrm{LeakyReLU}\big(a_r^{\top}[\,W_r h_i^{(l)};\,W_r h_j^{(l)}\,]\big)\Big),
\quad
h_i^{(l+1)}=\sigma\!\Big(\sum_{r\in\mathcal{R}}\sum_{j\in\mathcal{N}_i^{r}}\alpha_{ij}^{r}\,W_r h_j^{(l)}\Big).
\end{equation}
After $L$ layers we obtain contextualized representations $\tilde h_i\in\Real^{H}$,
written compactly as $\tilde H=\Enc(\{h_i^{0}\})$. A candidate-pair representation
concatenates the two utterance vectors with their element-wise product and absolute
difference,
\begin{equation}
h_{ij}^{\mathrm{pair}}=\big[\tilde h_i;\,\tilde h_j;\,\tilde h_i\odot\tilde h_j;\,|\tilde h_i-\tilde h_j|\big]\in\Real^{4H}.
\label{eq:pair}
\end{equation}

\paragraph{Self-contained pairs} For a self-cause pair ($j=i$) this representation
degenerates: the difference block $|\tilde h_i-\tilde h_i|$ is identically zero and the
product collapses to $\tilde h_i\odot\tilde h_i$, so a quarter of $h_{ii}^{\mathrm{pair}}$
carries no signal---an avoidable handicap given that such pairs are about half of the
labels. We therefore inject a learnable self-loop signature $e^{\mathrm{self}}\in\Real^{H}$
into the otherwise-dead difference block on the diagonal,
$h_{ii}^{\mathrm{pair}}=[\tilde h_i;\,\tilde h_i;\,\tilde h_i\odot\tilde h_i;\,e^{\mathrm{self}}]$,
which gives self-cause pairs a distinct marker while leaving the $4H$ dimension and
every cross-utterance pair unchanged. Initialized at $e^{\mathrm{self}}=0$, it recovers
the plain representation at the start of training, and the same signature is applied
identically in the factual and counterfactual scorings of Section~\ref{sec:necessity},
so the necessity difference is uncontaminated by the marker.

\subsection{Soft Candidate Gating}
\label{sec:gate}
Enumerating all pairs is quadratic and computing evidence for every pair is costly.
We define a soft relevance gate from the auxiliary heads,
$\pi_{ij}=\hat p^{\,e}_i\cdot \hat p^{\,c}_j\in[0,1]$, and allocate the evidence
computations to the top-$M$ pairs by $\pi_{ij}$, denoted $\mathcal{T}$. Gating governs
\emph{where computation is spent}, not which pairs may be positive: every valid pair
still contributes to the pair loss, so no candidate is permanently discarded.

\subsection{Perturbation-Based Necessity Evidence}
\label{sec:necessity}
A candidate cause should matter only if removing it changes the prediction. A scalar
pair-scoring head $f:\Real^{4H}\!\to\!\Real$ produces the factual score
$s_{ij}=f(h_{ij}^{\mathrm{pair}})$, supervised by the pair labels,
$\mathcal{L}_{\mathrm{score}}=\mathrm{BCE}(\sigma(s_{ij}),y_{ij})$.

\paragraph{Anchored baselines} Perturbations replace an element with a neutral
baseline rather than zero. With detached exponential-moving-average means $\bar h$
(utterance space) and $\bar x^{m}$ (modality-$m$ input space),
\begin{equation}
b^{\mathrm{utt}}=\bar h+\delta^{\mathrm{utt}},\quad
b^{m}=\bar x^{m}+\delta^{m},\quad
\mathcal{L}_{\mathrm{anc}}=\|\delta^{\mathrm{utt}}\|_2^2+\sum_{m}\|\delta^{m}\|_2^2,
\end{equation}
where the learnable residuals $\delta$ are penalized so the baselines remain near
genuine ``no-information'' means.

\paragraph{Cause necessity} Replacing the contextualized cause representation
$\tilde h_j$ by $b^{\mathrm{utt}}$ and re-scoring yields
\begin{equation}
s_{ij}^{-u}=f\!\left(\big[\tilde h_i;\,\tilde h_j^{-u};\,\tilde h_i\odot\tilde h_j^{-u};\,|\tilde h_i-\tilde h_j^{-u}|\big]\right),
\qquad
\Delta_{ij}^{u}=s_{ij}-s_{ij}^{-u}.
\end{equation}
Because the textual stream is the backbone of the encoder, this representation-level
term also captures the semantic necessity of the cause; we reserve modality
perturbation for the non-verbal streams.

\paragraph{Pair-local modality contribution} For $m\in\{a,v\}$ a faithful
perturbation must propagate through the non-linear encoder. We re-encode the whole
conversation once per modality with that modality replaced by its baseline at the
input, $\tilde H^{-m}=\Enc(\{h_i^{0}\}\mid x^{m}\!\leftarrow b^{m})$, costing two
extra encodings per conversation shared by all of its pairs. Keeping the original
emotion representation and using the perturbed cause representation $\tilde h_j^{-m}$,
\begin{equation}
\Delta_{ij}^{m}=s_{ij}-f\!\left(\big[\tilde h_i;\,\tilde h_j^{-m};\,\tilde h_i\odot\tilde h_j^{-m};\,|\tilde h_i-\tilde h_j^{-m}|\big]\right),
\end{equation}
which makes the modality term pair-local and directly comparable to the cause term.
The normalized terms and modality weights are
$\tilde\Delta_{ij}^{m}=(\Delta_{ij}^{m}-\mu_m)/(\sigma_m+\epsilon)$ and
$w_{ij}^{m}=\operatorname{softmax}_{m\in\{a,v\}}(\tilde\Delta_{ij}^{m})$, with running
statistics $(\mu_m,\sigma_m)$; the cause term is standardized within its own group as
$\tilde\Delta_{ij}^{u}=(\Delta_{ij}^{u}-\mu_u)/(\sigma_u+\epsilon)$.

\paragraph{Necessity-evidence vector} Collecting the terms gives
\begin{equation}
z_{ij}^{\mathrm{nec}}=\big[\,\tilde\Delta_{ij}^{u};\ \tilde\Delta_{ij}^{a};\ \tilde\Delta_{ij}^{v};\ w_{ij}^{a};\ w_{ij}^{v}\,\big]\in\Real^{5},
\end{equation}
consumed by the fusion module with a stop-gradient so that noisy early evidence
informs the classifier without destabilizing the backbone.

\paragraph{Necessity calibration} To ensure the necessity evidence is semantically
meaningful---removing a true cause decreases the score, removing a non-cause does
not---we add a hinge-style loss,
\begin{equation}
\mathcal{L}_{\mathrm{cal}}=\frac{1}{|\mathcal{B}|}\sum_{(i,j)\in\mathcal{B}}\Big[\mathbb{1}[y_{ij}=1]\,\mathrm{ReLU}(-\Delta_{ij}^{u})+\mathbb{1}[y_{ij}=0]\,\mathrm{ReLU}(\Delta_{ij}^{u}-\kappa_{\mathrm{cal}})\Big],
\end{equation}
with a small margin $\kappa_{\mathrm{cal}}\ge 0$. We stress that this necessity
evidence is a \emph{model-level} quantity: it measures the sensitivity of the model's
own prediction to controlled input perturbations, used to regularize and explain
predictions rather than to assert causal relations in the world.

\subsection{Bounded Positional Prior}
\label{sec:pos}
Causes tend to lie near the emotion utterance, so relative position is informative;
the risk is over-reliance. With normalized signed distance $d_{ij}=(i-j)/N\in[-1,1]$
and a small network $\psi$,
\begin{equation}
b_{ij}^{\mathrm{pos}}=\eta\cdot\tanh(\psi(d_{ij}))\in(-\eta,\eta),\qquad
\mathcal{L}_{\mathrm{pos}}=\frac{1}{|\mathcal{B}|}\sum_{(i,j)\in\mathcal{B}}\big(b_{ij}^{\mathrm{pos}}\big)^2 .
\end{equation}
The amplitude $\eta$ caps the prior and $\mathcal{L}_{\mathrm{pos}}$ discourages it
from dominating, so position can refine but not override the semantic evidence.

\subsection{LLM-Guided Counterfactual Evidence Distillation}
\label{sec:llm}
LLMs reason effectively over text but cannot perceive raw acoustic or visual signals;
we restrict their role to text-grounded reasoning and obtain modality evidence solely
from Section~\ref{sec:necessity}. Because querying an LLM at inference is impractical,
its reasoning is distilled offline into a lightweight student.

\paragraph{Counterfactual reasoning evidence} For a candidate pair, a fixed structured
prompt asks the LLM to reason counterfactually about an intervention on the candidate
cause and to return four calibrated scalars in $[0,1]$: necessity $s_{ij}^{\mathrm{nec}}$
(how likely the emotion in $u_i$ would disappear or weaken had $u_j$ not occurred),
sufficiency $s_{ij}^{\mathrm{suf}}$ (how likely $u_j$ alone elicits the emotion),
direction $s_{ij}^{\mathrm{dir}}$ (whether $u_j$ causes $u_i$ rather than the reverse),
and spuriousness $s_{ij}^{\mathrm{spur}}$ (probability that $u_i,u_j$ merely co-occur).
Reliability is the inter-sample agreement $\rho_{ij}\in[0,1]$ across $k$ stochastic
samples. The reasoning-evidence vector is
\begin{equation}
z_{ij}^{\mathrm{rea}}=\big[\,s_{ij}^{\mathrm{nec}};\ s_{ij}^{\mathrm{suf}};\ s_{ij}^{\mathrm{dir}};\ s_{ij}^{\mathrm{spur}};\ \rho_{ij}\,\big]\in[0,1]^{5}.
\end{equation}
These are \emph{textual} counterfactual judgements obtained by prompting, not a formal
structural-causal-model identification.

\paragraph{Self-contained causes} For self-cause pairs ($j=i$) the ``remove $u_j$''
framing is ill-posed, since deleting the utterance also deletes the emotion it carries.
We branch the prompt on whether $j=i$: for $j\neq i$ the four scalars retain the
cross-utterance reading; for $j=i$ the counterfactual is taken over the \emph{event
described} in $u_i$ rather than the sentence. Both branches share the same four-scalar
schema, so all downstream consumers are unchanged, and the self-pair $(i,i)$ is always
included among the annotation candidates for every non-neutral emotion utterance.

\paragraph{Offline annotation} Annotation is performed once with a frozen base model
$\theta_0$ (after warmup). A pair is \emph{hard} when the model is unconfident or its
decision conflicts with the necessity evidence,
\begin{equation}
\mathcal{P}_{\mathrm{hard}}=\big\{(i,j)\mid \max(p_{ij},1-p_{ij})<\tau \ \text{or}\ \mathrm{Conflict}_{ij}=1\big\},
\end{equation}
\begin{equation}
\mathrm{Conflict}_{ij}=\mathbb{1}\big[(p_{ij}\!\ge\!0.5\wedge \Delta_{ij}^{u}\!\le\!0)\vee(p_{ij}\!<\!0.5\wedge \Delta_{ij}^{u}\!>\!\kappa)\big].
\end{equation}
The annotated set is drawn from the population $\mathcal{T}$ seen at inference and
stratified into hard and easy pairs,
$\mathcal{S}=(\mathcal{T}\cap\mathcal{P}_{\mathrm{hard}})\cup\mathrm{sample}(\mathcal{T}\setminus\mathcal{P}_{\mathrm{hard}})$.

\paragraph{Distillation} A student head predicts the reasoning evidence from the pair
representation alone, $\hat z_{ij}^{\mathrm{rea}}=q(h_{ij}^{\mathrm{pair}})$, with
$q:\Real^{4H}\to[0,1]^{5}$, which prevents it from trivially copying the necessity
evidence. Distillation uses a reliability-weighted grouped binary cross-entropy,
\begin{equation}
\mathcal{L}_{\mathrm{dst}}=\frac{1}{|\mathcal{S}|}\sum_{(i,j)\in\mathcal{S}}\rho_{ij}\sum_{c=1}^{5}\mathrm{BCE}\!\left(\hat z_{ij,c}^{\mathrm{rea}},\,z_{ij,c}^{\mathrm{rea}}\right).
\end{equation}
At inference the LLM is never queried; the student supplies $\hat z_{ij}^{\mathrm{rea}}$,
whose fifth coordinate $\hat\rho_{ij}$ is an inference-time reliability estimate.

\subsection{Complementarity-Aware Evidence Fusion and Prediction}
\label{sec:fusion}
Fusion realizes an explicit complementarity: the LLM should compensate the small model
where the latter is uncertain, while the small model should lead where the LLM is
unreliable or disagrees with the model's own necessity evidence. For pairs in
$\mathcal{T}$ a smooth presence weight indicates how firmly a pair lies within the
evidence budget, $g^{\mathrm{ev}}_{ij}=\sigma(s_\pi(\pi_{ij}-\pi_M))$ with $a_{ij}=g^{\mathrm{ev}}_{ij}$,
and pairs outside $\mathcal{T}$ approach $a_{ij}\!\to\!0$.

\paragraph{Reliability scaling} The distilled evidence is attenuated by its predicted
reliability, scaling only the first four components and preserving $\hat\rho_{ij}$:
$\bar z_{ij}^{\mathrm{rea}}=[\hat\rho_{ij}\,\hat z_{ij,1:4}^{\mathrm{rea}};\,\hat\rho_{ij}]$.

\paragraph{Confidence and agreement} Reading the model's confidence from the factual
score $s_{ij}$,
\begin{equation}
c_{ij}=|2\sigma(s_{ij})-1|,\quad
\alpha_{ij}^{\mathrm{agr}}=1-|\sigma(\tilde\Delta_{ij}^{u})-s_{ij}^{\mathrm{nec}}|,\quad
\mathrm{cf}_{ij}=1-\alpha_{ij}^{\mathrm{agr}}.
\end{equation}
The fused evidence representation and the conditioned gate are
\begin{align}
h_{ij}^{\mathrm{ev}}&=\phi\!\left(\big[h_{ij}^{\mathrm{pair}};\,\sg(z_{ij}^{\mathrm{nec}});\,\bar z_{ij}^{\mathrm{rea}};\,a_{ij};\,c_{ij};\,\alpha_{ij}^{\mathrm{agr}}\big]\right),\\
g_{ij}&=\sigma\!\big(W_g[\,h_{ij}^{\mathrm{pair}};\,h_{ij}^{\mathrm{ev}};\,c_{ij};\,\hat\rho_{ij};\,\mathrm{cf}_{ij}\,]+b_g\big),\\
h_{ij}^{\mathrm{final}}&=g_{ij}\odot h_{ij}^{\mathrm{pair}}+(1-g_{ij})\odot h_{ij}^{\mathrm{ev}},
\end{align}
so a less confident small model can defer to the evidence, whereas an unreliable or
conflicting LLM is suppressed in favor of the pair representation.

\paragraph{Necessity-modulated agreement} A naive agreement boost is direction-agnostic;
we modulate it by the average necessity of the two sources,
$\bar n_{ij}=\tfrac{1}{2}(\sigma(\tilde\Delta_{ij}^{u})+s_{ij}^{\mathrm{nec}})$, and define
\begin{equation}
p_{ij}=\sigma\!\big(w_f^{\top}h_{ij}^{\mathrm{final}}+b_f+b_{ij}^{\mathrm{pos}}+w_\alpha\,(2\alpha_{ij}^{\mathrm{agr}}-1)\,\bar n_{ij}\big),
\end{equation}
with a scalar $w_\alpha$ initialized at $0$. When both sources agree the pair is causal
($\bar n_{ij}\!\approx\!1$) the boost is active; when both agree it is not
($\bar n_{ij}\!\approx\!0$) the boost vanishes, preventing false-positive inflation.
Dropping the conditioning signals recovers an unconditioned gate, which we retain as an
ablation baseline.

\subsection{Training and Inference}
\label{sec:training}
The pair objective is a class-weighted BCE over all positives and a sampled set of
negatives at ratio $r$ within the valid upper-triangular region,
$\mathcal{L}_{\mathrm{pair}}=\frac{1}{|\mathcal{P}^{+}\cup\mathcal{N}_r|}\sum \mathrm{wBCE}_{\omega}(p_{ij},y_{ij})$.
With auxiliary emotion/cause objectives and the preliminary emotion head,
$\mathcal{L}_{\mathrm{cls}}=\mathcal{L}_{\mathrm{pair}}+\alpha_1\mathcal{L}_{\mathrm{emo}}+\alpha_2\mathcal{L}_{\mathrm{cau}}+\mathcal{L}_{\mathrm{emo}}^{\mathrm{pre}}$,
and the overall objective is
\begin{equation}
\mathcal{L}=\mathcal{L}_{\mathrm{cls}}+\beta\,\mathcal{L}_{\mathrm{score}}+\lambda_1\,\mathcal{L}_{\mathrm{dst}}+\lambda_2\,\mathcal{L}_{\mathrm{pos}}+\gamma\,\mathcal{L}_{\mathrm{anc}}+\gamma\,\mathcal{L}_{\mathrm{cal}}.
\end{equation}
Training proceeds in two phases. During the first $K$ warmup epochs the encoder and the
auxiliary/scoring heads are trained with $\mathcal{L}_{\mathrm{cls}}$ and
$\mathcal{L}_{\mathrm{score}}$, the graph uses only structural relations, and evidence
fusion is disabled; the student is nonetheless trained from epoch~1 via
$\mathcal{L}_{\mathrm{dst}}$. The base model is then frozen to produce the offline LLM
annotations. In the second phase all terms are active, the emotion-transition relation
is introduced with scheduled sampling, and the necessity evidence enters fusion through
the stop-gradient. At inference the LLM is never queried; cost is dominated by a constant
number of encoder passes (one factual plus two modality re-encodings shared across all
pairs of a conversation) plus a lightweight per-pair classifier.

% =====================================================================
\section{Experiments}
\label{sec:exp}

\noindent\textbf{Note (draft).} All numbers in this section are placeholders to be
replaced with the final experimental results; they are reported as
mean$\pm$standard-deviation over three random seeds. \todo{replace every number in
Tables~\ref{tab:main}--\ref{tab:analysis} with real results from
\texttt{scripts/aggregate\_results.py}}

\subsection{Dataset and Evaluation}
We evaluate on the benchmark Emotion-Cause-in-Friends (ECF) corpus
\citep{wang2023mecpe}, which annotates emotion-cause pairs over multi-party
conversations from the TV series \emph{Friends}, with aligned textual, acoustic, and
visual signals. We follow the standard split (1{,}001 training conversations and the
official development/test sets). We report precision, recall, and $F_1$ for the
emotion-cause pair extraction task (\textbf{Pair}), as well as auxiliary emotion-utterance
$F_1$ (\textbf{Emo}) and cause-utterance $F_1$ (\textbf{Cause}). Following best
practice, the decision threshold is tuned on the development set and then fixed on the
test set, avoiding in-sample threshold leakage; all numbers are averaged over three
seeds.

\subsection{Implementation Details}
The text encoder is BERT-base (cased) \citep{devlin2019bert}. Acoustic features are
$6373$-dimensional openSMILE descriptors \citep{eyben2010opensmile} and visual features
are $4096$-dimensional 3D-CNN features \citep{tran2015c3d}. The hidden size is $H=400$
and the relation-aware encoder uses multi-head attention. We optimize with AdamW
\citep{loshchilov2019adamw}, a BERT learning rate of $1\!\times\!10^{-5}$, a head
learning rate of $1\!\times\!10^{-4}$, batch size $4$, and $25$ epochs. The four exposed
hyperparameters are the warmup length $K\!=\!3$, the evidence budget $M\!=\!50$, the
distillation weight $\lambda_1\!=\!1$, and the fusion mode (complementarity-aware).
Remaining coefficients are fixed at grounded defaults
($\alpha_1\!=\!\alpha_2\!=\!\beta\!=\!1$, $\lambda_2\!=\!0.1$, $\gamma\!=\!0.01$,
$\eta\!=\!1$, negative ratio $r\!=\!5$, class weight $\omega\!=\!2$). Offline LLM
annotations are produced once with $k$ stochastic samples per hard pair. \todo{confirm
$k$ and the exact LLM used}

\subsection{Main Results}
Table~\ref{tab:main} compares the proposed method with a legacy graph baseline and a
published baseline. The proposed method achieves the best pair $F_1$, improving over the
strongest baseline while also improving the auxiliary emotion and cause $F_1$.
% TODO: replace placeholders below with real numbers
\begin{table}[t]
\centering
\caption{Main results on the ECF test set (\%, mean over three seeds). All values are
placeholders. \todo{fill from \texttt{results/TABLES.md} Table~1}}
\label{tab:main}
\begin{tabular}{lccccc}
\toprule
Method & Pair P & Pair R & Pair $F_1$ & Emo $F_1$ & Cause $F_1$ \\
\midrule
Published baseline \citep{li2024hilo} & 53.10 & 56.20 & 54.61 & 69.80 & 67.90 \\
Graph (legacy, \texttt{use\_method=no}) & 53.80 & 56.90 & 55.31 & 70.10 & 68.20 \\
\textbf{Ours (full method)} & \textbf{55.40} & \textbf{58.30} & \textbf{56.81} & \textbf{71.20} & \textbf{69.40} \\
\bottomrule
\end{tabular}
\end{table}

\subsection{Ablation Study}
Table~\ref{tab:ablation} removes one component at a time from the full model. Every
component contributes positively to the pair $F_1$; the necessity evidence and the
distilled reasoning evidence yield the largest drops when removed, and replacing the
complementarity-aware fusion with the unconditioned gate also degrades performance.
% TODO: replace placeholders with real ablation numbers
\begin{table}[t]
\centering
\caption{Ablation on the ECF test set (\%). $\Delta$ is the change in pair $F_1$ relative
to the full model. All values are placeholders. \todo{fill from
\texttt{results/ablation\_table.md}}}
\label{tab:ablation}
\begin{tabular}{lccccc}
\toprule
Configuration & Pair P & Pair R & Pair $F_1$ & Cause $F_1$ & $\Delta$ Pair $F_1$ \\
\midrule
\textbf{Full} & 55.40 & 58.30 & 56.81 & 69.40 & --- \\
\quad $-$ Self-loop signature (\S\ref{sec:encoder}) & 55.00 & 57.80 & 56.37 & 69.00 & $-0.44$ \\
\quad $-$ Emotion transition (\S\ref{sec:encoder}) & 54.80 & 57.70 & 56.21 & 68.90 & $-0.60$ \\
\quad $-$ Necessity evidence (\S\ref{sec:necessity}) & 54.30 & 57.10 & 55.66 & 68.40 & $-1.15$ \\
\quad $-$ Positional prior (\S\ref{sec:pos}) & 55.10 & 57.90 & 56.46 & 69.10 & $-0.35$ \\
\quad $-$ Distillation (\S\ref{sec:llm}) & 54.10 & 57.30 & 55.65 & 68.30 & $-1.16$ \\
\quad Fusion $\rightarrow$ uncond.\ gate (\S\ref{sec:fusion}) & 54.70 & 57.60 & 56.11 & 68.80 & $-0.70$ \\
\bottomrule
\end{tabular}
\end{table}

\subsection{Analysis}
\paragraph{Self-contained vs.\ cross-utterance causes}
Table~\ref{tab:analysis} breaks down pair $F_1$ by whether the cause is self-contained
($j=i$) or cross-utterance ($j\neq i$). The self-loop signature and the branched
counterfactual prompt jointly improve the self-contained subset, which constitutes about
half of the gold labels. \todo{fill the breakdown}
\begin{table}[t]
\centering
\caption{Pair $F_1$ by cause type on ECF (\%). Placeholders. \todo{fill}}
\label{tab:analysis}
\begin{tabular}{lccc}
\toprule
Method & Self-cause $F_1$ & Cross-cause $F_1$ & Overall $F_1$ \\
\midrule
Graph (legacy) & 58.90 & 51.20 & 55.31 \\
\textbf{Ours} & \textbf{60.80} & \textbf{52.40} & \textbf{56.81} \\
\bottomrule
\end{tabular}
\end{table}

\paragraph{Effect of teacher reliability}
Distillation weighted by the LLM self-agreement $\rho_{ij}$ improves the net benefit of
the reasoning evidence relative to unweighted distillation, indicating that down-weighting
unreliable teacher judgements helps. \todo{add the reliability-vs-gain plot/table}

\paragraph{Qualitative example}
\todo{add one or two qualitative cases showing where the necessity evidence and the LLM
counterfactual judgement complement each other.}

% =====================================================================
\section{Conclusion}
\label{sec:conclusion}
We presented a multimodal emotion-cause pair extraction model that grounds its pair
decisions in three complementary sources of evidence: model-internal necessity evidence
obtained by faithful, pair-local perturbations; a bounded positional prior; and
text-grounded counterfactual reasoning distilled offline from a large language model. A
complementarity-aware gating network fuses these signals under explicit confidence,
reliability, and agreement conditioning, so that each source compensates the others'
weaknesses, while the LLM is never queried at inference. Experiments on the ECF benchmark
show consistent improvements over strong baselines, and ablations confirm that every
component contributes. % TODO: finalize the quantitative summary
In future work we will study the interaction between the distilled reasoning and stronger
multimodal encoders, and extend the necessity evidence to multi-hop causal chains.

% =====================================================================
\section*{Acknowledgments}
\todo{add acknowledgments and funding information if applicable.}

\bibliographystyle{elsarticle-harv}
\bibliography{refs}

\end{document}
